\begin{abstract}
Vision-language models (VLMs) are increasingly deployed as autonomous agents that interact with graphical user interfaces, yet their robustness to the natural visual variability of real-world UIs remains poorly understood. We introduce \bench{}, a controlled micro-benchmark comprising 300 UI screenshots (50 base pages across five categories, each rendered under five calibrated drift variants) and 390 grounded question-answer pairs. Each drift type, spanning theme changes, sidebar toggles, responsive scaling, viewport crops, and composite perturbations, is assigned a severity level from one to five, enabling fine-grained analysis of degradation curves. We evaluate three frontier VLMs (GPT-4o, GPT-4o-mini, Gemini~2.0~Flash) on answer accuracy, answer consistency, and evidence grounding IoU. Our experiments reveal that drift robustness does not correlate perfectly with base accuracy: Gemini~2.0~Flash achieves the highest Drift Robustness Score (0.950) despite lower base accuracy than GPT-4o. Critically, evidence grounding IoU remains below 0.07 across all models and conditions, exposing a fundamental gap in spatial evidence localization that must be addressed before VLM agents can be trusted in audit-critical workflows.
\end{abstract}
